\section{Multivariable Calculus}

Here, we include no proofs.

\begin{exercise}
    Find the equation of the plane in $\R^3$ passing through 
    $P=(0,0,1), Q=(1,0,0), R=(1,1,1)$.
\end{exercise}
\begin{solution}
    We use a reference point and use the vector cross product:
    \[
    v := \overrightarrow{PQ} = \langle 1,0,-1\rangle, \qquad 
    w := \overrightarrow{PR} = \langle 1, 1, 0\rangle
    \]
    Now, the normal vector to $P$ is $v \times w$:
    \[
    v \times w = \langle 1,-1,1 \rangle
    \]
    The equation of the plane is given by 
    \[
    \langle 1,-1,1\rangle \cdot \langle x-0, y-0, z-1 \rangle,
    \]
    which is simply $x-y+z = 1$.
\end{solution}

\begin{exercise}
    Compute the arc length of the helix $\gamma(t) = (\sin(t),-\cos(t), t)$
    from $t = 0$ to $t = 2\pi$.
\end{exercise}
\begin{solution}
    We must use the arc-length formula $\int_t \lVert \gamma'(t) \rVert dt$.
    \[
    \int_{t = 0}^{2\pi} \sqrt{\cos^2(t) + \sin^2(t) + 1^2} dt 
        = \int_{t = 0}^{2\pi} \sqrt{2} dt = 2\pi \sqrt{2}.
    \]
\end{solution}

The following is a nice way to determine if a limit doesn't exist:
\begin{theorem}
    $\lim \limits_{(x,y) \rightarrow (a,b)} f(x,y) = L$ if and only if for all
    $\gamma \colon [-,1,1] \rightarrow \R^2$ such that $\gamma(0) = (a,b)$, then 
    $\lim \limits_{t \to 0} f(\gamma(t)) = L$.
\end{theorem}

\begin{exercise}
Prove that $\lim \limits_{(x,y) \rightarrow (0,0)} \frac{xy}{x^2 + y^2}$
does not exist.
\end{exercise}
\begin{solution}
We consider the paths $\gamma(t) = (t,t)$ and $\delta(t) = (t,-t)$, which both 
tend to $(0,0)$ as $t \rightarrow 0$. Using $\gamma$, we have as follows:
\[
\lim \limits_{(x,y) \rightarrow (0,0)} \frac{xy}{x^2 + y^2} =
\lim \limits_{t \rightarrow 0} \frac{t^2}{t^2 + t^2} = \frac{1}{2}
\]    
On the other hand, $\delta$ gives:
\[
\lim \limits_{(x,y) \rightarrow (0,0)} \frac{xy}{x^2 + y^2} =
\lim \limits_{t \rightarrow 0} \frac{-t^2}{t^2 + t^2} = -\frac{1}{2}
\]    
Therefore, the limit doesn't exist.
\end{solution}

\begin{exercise}
Find $\lim \limits_{(x,y) \rightarrow (0,0)} \frac{xy^2}{x^2 + y^2}$.
\end{exercise}
\begin{solution}
We use polar coordinates and change the limit to $r \to 0$.
\[
\lim \limits_{(x,y) \rightarrow (0,0)} \frac{xy^2}{x^2 + y^2} = 
\lim \limits_{r \rightarrow 0} \frac{r^3\cos \theta \sin^2\theta}{r^2} =
\lim \limits_{r \rightarrow 0} r\cos \theta \sin^2 \theta
\]
Now,
\[
0 \leq |\lim \limits_{r \rightarrow 0} r\cos \theta \sin^2 \theta| \leq 
|\lim \limits_{r \rightarrow 0} r | = 0,
\]
so the limit is $0$ via squeeze.
\end{solution}

\subsection{Grad, Div, Curl}

Recall that the gradient satisfies the following:

\begin{proposition}
\[
\nabla \phi = \left\langle\frac{\partial \phi}{\partial x},
\frac{\partial \phi}{\partial y},  \frac{\partial \phi}{\partial z} \right\rangle
\]
Further, if we want the directional derivative of $\phi$ in the deriction of a
unit vector $\hat{s}$, denoted $\pfrac{\phi}{s}$, we have
\[
\pfrac{\phi}{s} = \hat{s} \cdot \nabla \phi.
\]
\end{proposition}


\begin{exercise}
    If $\phi = axy^2 + byz + cx^3z^2$, where $a,b,c$ are constants, find $\nabla \phi$. 
    Further, find the directional derivative of $\phi$ at the point $(1,4,2)$ in the 
    direction towards the point $(2,0-1)$.
\end{exercise}
\begin{solution}
    First, we compute the gradient:
    \[
    \nabla \phi = \left\langle ay^2 + 3cz^2z^2, 2axy + bz, by + 2cx^3z \right\rangle
    \]
    Therefore, the gradient at $P = (1,4,2)$ is:
    \[
    (\nabla \phi)_P = \left\langle 16a + 12c, 8a + 2b, 4b + 4c \right\rangle
    \]
    Now, $s$ is the difference between the two points, which is just
     $s = (2,0,-1) - (1,4,2) = (1,-4,-3)$. Thus, $\hat{s}$ is
    \[
    \hat{s} = \frac{\langle 1, -4 , -3\rangle}{\sqrt{26}}
    \]
    Finally, the directional derivative $\frac{\partial \phi}{\partial s}$, is as follows:
    \[
    \frac{\partial \phi}{\partial s} = (\nabla \phi \cdot \hat{s})_P = \frac{(16a + 12c) - 4(8a + 2b) - 3 (4b + 4c)}{\sqrt{26}} = \frac{-16a - 20b}{\sqrt{26}}
    \]
\end{solution}

\begin{df}
    Let $F$ be a vector field.
    \begin{itemize}
        \item The \emph{divergence} of $F$ is the scalar product 
        $\nabla \cdot F$, writing 
        $\divergence F = \nabla \cdot F$.
        \item The \emph{curl} of $F$ is the vector product 
        $\nabla \times F$, writing $\curl F = \nabla \times F$.
    \end{itemize}
\end{df}

We recall some useful identities relating these different notions.

\begin{theorem}
    Let $\phi \colon \R^3 \rightarrow \R$ and $F$ be a vector field (or 
    more simply, a vector valued function).
    \begin{itemize}
        \item $\curl (\nabla \phi) = 0$
        \item $\divergence (\curl F) = 0$
        \item $\curl (\curl (F)) = \nabla (\divergence F) - \nabla ^2 F$
    \end{itemize}
\end{theorem}